%% System design — data, agent (neutron-current focus), then RAG stubs
\subsection{Data description}
\label{sec:data}
Two complementary resources are used throughout this study: plant neutron-current measurements for
agent-side diagnosis, and a curated nuclear I\&C manual corpus for manual-grounded RAG and retrieval
evaluation.

\paragraph{Neutron-current measurement archive.}
The underlying measurement archive is organized into four groups of multi-channel neutron-current CSV
records (hereafter Groups~A1, A2, B1, and B2), corresponding to different instrumentation banks in the
same plant campaign.
In total, the archive contains 42 series files, as summarized in~\cref{tab:neutron-data}.

\begin{table}[H]
  \centering
  \caption{Composition of the neutron-current measurement archive.}
  \label{tab:neutron-data}
  \begin{tabular}{@{}lccc@{}}
    \toprule
    Group & Number of series & Channels per series & Columns per series \\
    \midrule
    A1 & 11 & 7 & 8 \\
    A2 & 11 & 7 & 8 \\
    B1 & 10 & 7 & 8 \\
    B2 & 10 & 7 & 8 \\
    \midrule
    Total & 42 & --- & --- \\
    \bottomrule
  \end{tabular}
\end{table}

Each series stores a time column together with seven neutron-current channels (eight columns overall);
channel identifiers follow plant tagging conventions (e.g., \texttt{RIC}\ldots\texttt{MA\_VER}).
Within each group, the seven detectors are arranged in a circular (circumferential) layout around the
reactor, so neighboring channels correspond to adjacent angular positions rather than to an arbitrary
listing order in the CSV.
Representative series contain on the order of $6\times 10^{5}$ samples at a nominal interval of
\SI{1}{\second} and cover multi-day windows.

Importantly, timestamps in the \emph{raw} CSV files are only approximately chronological: local
disorders and non-monotonic segments appear, so the records are not strictly sorted in increasing time
order as exported from the plant historian.
The agent therefore begins with an explicit chronological sort before any downstream probe is applied.

Plant-side operational annotations and fault-related records exist for parts of this campaign; however,
they are heterogeneous across sources and are not aligned into a uniform, sample-level fault taxonomy
that could be tabulated without substantial manual curation.
We therefore do \emph{not} report fault-label counts in this paper.
Instead, the zero-prior agent treats channel tags as opaque identifiers and does not consume fault labels
during analysis; any available annotations are reserved for future supervised benchmarking rather than
for the present pipeline.

For the system demonstration and case studies in~\cref{sec:results}, we focus on a working subset drawn
from Groups~A1 and A2, processed independently in workspace-local analyses.
The same agent protocol applies unchanged to the remaining series in the archive.

\paragraph{Nuclear I\&C manual corpus.}
The knowledge corpus comprises 59 technical documents converted from PDF to Markdown for indexing
and is partitioned into a \emph{core} pool and a \emph{non-core} (supplementary) pool.
The core pool contains 13 IAEA manuals, including the IAEA OLM reports on instrument-channel
monitoring and on process/component condition monitoring and
diagnostics~\citep{iaea2008olm1,iaea2008olm2}, together with related IAEA guidance on I\&C design,
digital modernization, maintenance/surveillance, and operational limits and procedures.
The non-core pool contains the remaining 46 same-domain manuals drawn from NRC, EPRI/DOE,
OECD/NEA, and related laboratory or topical materials on sensors, noise, and core/instrument
monitoring.
These non-core documents are \emph{not} treated merely as noise to be rejected; rather, they are
topically related supplements that enlarge the searchable knowledge base in a controlled way.
A central experimental question of this work is therefore the \emph{core--non-core relationship}:
how retrieval behavior changes as the ratio of core to non-core documents varies---equivalently,
how large a non-core supplement can be admitted for a fixed core set before target (core) evidence
is diluted.
This controlled expansion underpins the corpus-scope and self-retrieval experiments
in~\cref{sec:eval}.

\subsection{Zero-prior diagnostic agent}
\label{sec:agent}
\subsubsection{Design objective}
The measurement-side goal is to turn a raw neutron-current CSV into a structured, human-readable
condition report without assuming a labeled fault library, a fixed channel-to-physics map, or
handcrafted frequency bands.
Channel tags are treated as opaque identifiers; statistical morphology, redundancy, anomalies, and
spatiotemporal cues are discovered from the series themselves.
This ``zero-prior'' stance matches the archive described in~\cref{sec:data}, where fault annotations
are heterogeneous and are deliberately not consumed at inference time.

\subsubsection{Architecture and control loop}
The agent couples a local LLM planner with a deterministic tool router~\citep{yao2023react}.
The planner proposes tool calls; a Python router executes registered analyzers and returns concise
status messages that update a persistent workspace blackboard (\texttt{state.json}).
Only the entry probe may receive a raw file path; thereafter every probe must take a workspace
directory alone and resolve the active CSV from the blackboard.
\Cref{fig:agent-workflow} summarizes the control flow.

After initialization, the planner is not allowed to restate free-text file paths.
LLM tool use remains brittle: models may invent nonexistent APIs, pass malformed arguments, or
generate fluent but unfaithful content~\citep{patil2024gorilla,schick2023toolformer,ji2023hallucination}.
The workspace contract therefore keeps authoritative paths and trusted state on the runtime side,
while the planner selects only the next registered probe; the failure modes that motivated these
constraints are detailed in~\cref{sec:agent-design}.

\begin{figure}[H]
  \centering
  \includegraphics[width=0.92\linewidth]{fig2_agent_workflow}
  \caption{Zero-prior agent workflow for neutron-current CSV analysis. Profiling creates a workspace
  blackboard; later probes operate only on the workspace directory and accumulate a condition report.}
  \label{fig:agent-workflow}
\end{figure}

\subsubsection{Workspace initialization}
Given a CSV path, the load-and-profile probe creates a dedicated workspace directory, parses the time
column, sorts the series into increasing chronological order (addressing the raw non-monotonic exports
noted in~\cref{sec:data}), and writes a cleaned active file.
It also records basic health indicators---shape, nominal sampling interval, missingness, and inactive
channels---on the blackboard and appends an initialization block to the condition report.
Subsequent probes refuse to run if the blackboard is missing, which enforces a mandatory warm-up step.

\subsubsection{Statistical discovery and anomaly probes}
After initialization, the planner may invoke the probes listed in~\cref{tab:agent-tools}.
The zero-prior discovery probe performs a global longitudinal scan: uniqueness and coefficient of
variation are used to assign coarse morphological roles (discrete/step-like, high-inertia, highly
volatile, or ordinary continuous), while pairwise correlations flag near-redundant channel pairs
($|r|\ge 0.99$) and strongly co-moving groups.
Given the circular detector layout, such collinearity often indicates angularly adjacent or
functionally redundant sensors rather than coincidental numerical agreement.

The anomaly/regime probe then segments each channel into steady and transient regimes and enumerates
spike episodes and isolated zero dropouts, summarizing which channels dominate each anomaly type.
The time-lag scan computes cross-correlations across channels to separate lead--lag conduction links
from essentially synchronous groups.
The extremes scan localizes global maxima/minima and their timestamps; the spatial-snapshot probe
freezes a multi-channel cross-section at a chosen time, which is useful after an extreme has been
identified.
Together, these probes accumulate evidence that can later be verbalized as spikes, dropouts,
redundancy, and regime changes without supervised class labels.

\begin{table}[H]
  \centering
  \footnotesize
  \setlength{\tabcolsep}{3pt}
  \caption{Diagnostic probes used by the neutron-current agent. The last column records the registered
  implementation identifier for reproducibility.}
  \label{tab:agent-tools}
  \begin{tabular}{@{}>{\raggedright\arraybackslash}p{0.19\linewidth}
                    >{\raggedright\arraybackslash}p{0.40\linewidth}
                    >{\raggedright\arraybackslash\ttfamily\scriptsize}p{0.33\linewidth}@{}}
    \toprule
    \textrm{\normalfont Probe} &
    \textrm{\normalfont Description} &
    \textrm{\normalfont Implementation ID} \\
    \midrule
    Load-and-profile &
    Creates workspace, sorts timestamps, records data-health indicators. &
    \path{load_and_profile_csv_tool} \\
    Zero-prior discovery &
    Infers morphological roles and co-moving or near-redundant groups. &
    \path{zero_prior_discovery_tool} \\
    Anomaly/regime &
    Segments regimes; enumerates spikes and zero dropouts. &
    \path{anomaly_detector} \\
    Time-lag scan &
    Estimates lead--lag versus synchronous groups. &
    \path{global_time_lag_scan_tool} \\
    Extremes scan &
    Localizes global maxima/minima and timestamps. &
    \path{global_extremes_scan_tool} \\
    Spatial snapshot &
    Multi-channel cross-section at a selected time. &
    \path{get_spatial_snapshot_tool} \\
    \bottomrule
  \end{tabular}
\end{table}

\subsubsection{Design process and adopted constraints}
\label{sec:agent-design}
The architecture in~\cref{sec:agent} is the outcome of an iterative, failure-driven design process
rather than a single a~priori specification.
Early prototypes delegated almost all bookkeeping to the planner: each probe call restated an absolute
CSV path, tool replies were expected to be nearly pure JSON, and the resulting condition text was
passed almost unchanged into downstream diagnosis.
Under a local mid-size LLM these assumptions failed repeatedly on multi-step neutron-current analyses.
\Cref{tab:agent-design} summarizes the issues retained as hard constraints in the final protocol;
the remainder of this subsection elaborates the failure mechanisms that motivate those rows.

\begin{table}[H]
  \centering
  \footnotesize
  \caption{Design issues encountered in early agent prototypes and constraints adopted in the final
  measurement-side protocol.}
  \label{tab:agent-design}
  \begin{tabular}{@{}p{0.30\linewidth}p{0.62\linewidth}@{}}
    \toprule
    Design issue & Adopted constraint \\
    \midrule
    Path restatement errors across tool steps
    (wrong folder; raw vs.\ sorted copy) &
    Only load-and-profile accepts a CSV path; later probes take a workspace directory and resolve the
    active file from the blackboard. \\
    Dropped or malformed tool JSON from local LLMs &
    Strip Markdown fences; extract the first brace-balanced JSON object/array; accept only registered
    probe identifiers. \\
    Oversized condition text for downstream reasoning &
    Retain a human-readable report; drive RAG from a compact structured alert summary. \\
    Ungrounded diagnostic prose without manuals &
    Build queries from the structured summary; require multi-stage, citation-bearing retrieval over the
    I\&C corpus. \\
    \bottomrule
  \end{tabular}
\end{table}

When every tool step re-emitted a free-text path, small transcription errors compounded over a probe
chain: the planner confused campaign folders, mixed raw exports with chronologically sorted copies, or
pointed later probes at a different series than the one just profiled.
Because the runtime trusted the path argument, such mistakes executed silently and produced incoherent
reports rather than explicit refusals.
The adopted remedy is to treat the load-and-profile step as the sole path-accepting entry point.
It creates a dedicated workspace, writes the cleaned active file, and records the authoritative
location on the blackboard; thereafter probes accept only a workspace directory and resolve the series
from that board~\citep{yao2023react,patil2024gorilla}.
Path ownership thus moves from mutable model strings to a runtime contract, which also makes incomplete
warm-up fail closed when the blackboard is missing.

A second class of failures arose at the action interface.
Strict ``JSON-only'' prompting was brittle: local models frequently wrapped tool calls in Markdown
fences, interleaved brief natural-language commentary, or truncated braces, so the router discarded
otherwise valid actions; unrestricted free text made parsing equally unreliable.
The present extractor therefore strips fences, searches for the first brace-balanced JSON object or
array, and validates that the \texttt{action} field matches a registered probe identifier before
execution~\citep{yao2023react,schick2023toolformer}.
Invented analyzer names fail closed.
This does not claim perfect structured generation; it treats imperfect local-model output as the
default case and recovers the call when a well-formed object is present.

Even when the probe chain completed, an unfiltered multi-page condition narrative overloaded the
manual-diagnosis stage: the LLM was asked to reason over lengthy prose while simultaneously selecting
handbook passages, which encouraged vague, poorly grounded summaries.
The design therefore separates audiences.
The workspace retains a chronological, human-readable report for inspection and audit, while a compact
structured summary---variable-wise alerts, spike/zero tallies, and collinearity or lag cues---is
reserved as the machine interface to RAG (\cref{sec:rag-diag}).
Retrieval queries are derived from that summary rather than from an operator's free-text guess about
which manual clause applies, and the synthesizer is required to attach numbered citations to retrieved
chunks~\citep{ji2023hallucination}.
Illustrative workspaces retained for this study show complete probe chains on representative A1/A2
series, while other series exhibit partial completion or redundant re-invocation of the same probe;
these cases indicate that the workspace and JSON contracts are enforceable, yet planner discipline
still affects coverage---a qualitative observation that motivates the protocol, not a supervised
accuracy claim with fabricated recovery rates.

\subsubsection{Calling protocol and outputs}
Tool calls must appear as constrained JSON objects
(\texttt{\{"action": \ldots, "parameters": \ldots\}}), optionally wrapped in Markdown fences; the
runtime extracts the first balanced JSON object or array even when the LLM adds free-text commentary,
mitigating the JSON-loss issue discussed above.
Only registered probe identifiers (\cref{tab:agent-tools}) are accepted by the local router
(Ollama~\texttt{qwen2.5:7b} in our deployment), so invented analyzer names fail closed rather than
executing an undefined path.
Each successful probe appends narrative blocks to the workspace report and updates structured fields on
the blackboard (e.g., morphological roles, spike/zero summaries, lag groups).
The resulting condition report---together with optional JSON snapshots---forms the sole
measurement-side interface to the RAG stage in~\cref{sec:rag-diag}: downstream retrieval is driven by
observed phenomena rather than by an operator's free-text guess about which manual clause applies.

\subsection{Manual corpus and index construction}
\label{sec:index}
Source PDFs are page-cropped when needed, converted to Markdown, chunked with a size/overlap grid, and
embedded with \texttt{nomic-embed-text} into a LlamaIndex vector store~\citep{liu2023llamaindex}.
Indexing is applied to the corpus described in~\cref{sec:data}; evaluation varies which documents are
admitted into the active retrieval pool while keeping the embedding model fixed.

\subsection{Report-driven multi-stage RAG diagnosis}
\label{sec:rag-diag}
Condition reports are parsed into compact JSON (variable-wise alerts, spike summaries,
topology/collinearity cues).
Queries are constructed from this JSON rather than from short free-text alone.
Retrieval proceeds in stages (phenomenon/mechanism passages, then redundancy/topology supplements)
and the LLM composes a multi-part diagnosis with numbered citations (\texttt{[n]}, \texttt{Rn}) linked
to retrieved chunks.
Because diagnosis quality depends on retrieval, \cref{sec:eval} evaluates the knowledge layer
separately.
